\documentclass[10pt]{article}
\usepackage[paperwidth=384pt,paperheight=900pt,margin=10pt]{geometry}
\usepackage[T1]{fontenc}
\usepackage{lmodern}
\usepackage[utf8]{inputenc}
\usepackage{amsmath,amssymb}
\usepackage{array}
\usepackage{enumitem}
\usepackage{xcolor}
\usepackage{microtype}

\setlength{\parindent}{0pt}
\setlength{\parskip}{4pt}
\setlist[itemize]{leftmargin=12pt}
\setlist[enumerate]{leftmargin=14pt}

\sloppy
\allowdisplaybreaks
\emergencystretch=2em

\begin{document}

\section*{Uppgift 1}

\subsection*{Original question}
Låt $A$ och $B$ vara två händelser sådana att
$P(A^c\mid B^c)=0.6$, $P(B^c\mid A^c)=0.4$ och
$P(A^c\cup B^c)=0.8$.
Beräkna $P(A)$ och $P(B)$.
\hfill (6p)

\subsection*{Compiled notes from Yacoub + Obid}

\subsubsection*{What we are doing}
We are given conditional probabilities with complements:
$A^c$ and $B^c$.

The notes first rewrite the conditional probabilities using
the formula

\[
P(E\mid F)
=\frac{P(E\cap F)}{P(F)}.
\]

Then we use that

\[
A^c\cap B^c
=
B^c\cap A^c
\]

is the same intersection.

Finally we use the union formula:

\[
P(A^c\cup B^c)
=
P(A^c)+P(B^c)-P(A^c\cap B^c).
\]

\subsubsection*{Given information}
From the question:

\[
P(A^c\mid B^c)=0.6
\]

\[
P(B^c\mid A^c)=0.4
\]

\[
P(A^c\cup B^c)=0.8
\]

We want:

\[
P(A)\quad \text{and}\quad P(B).
\]

\subsubsection*{Step-by-step solution}

\textbf{1. Rewrite the first conditional probability}

Using the conditional probability formula:

\begin{align*}
P(A^c\mid B^c)
&=
\frac{P(A^c\cap B^c)}
{P(B^c)} \\[2mm]
&=0.6
\end{align*}

Therefore:

\begin{align*}
P(A^c\cap B^c)
&=
0.6P(B^c)
\end{align*}

\textbf{2. Rewrite the second conditional probability}

\begin{align*}
P(B^c\mid A^c)
&=
\frac{P(B^c\cap A^c)}
{P(A^c)} \\[2mm]
&=0.4
\end{align*}

Therefore:

\begin{align*}
P(B^c\cap A^c)
&=
0.4P(A^c)
\end{align*}

But this is the same intersection:

\[
P(A^c\cap B^c)
=
P(B^c\cap A^c).
\]

So:

\begin{align*}
0.6P(B^c)
&=
0.4P(A^c)
\end{align*}

\textbf{3. Isolate one complement probability}

The notes isolate it in two useful ways.

First:

\begin{align*}
P(B^c)
&=
\frac{0.4}{0.6}P(A^c) \\[2mm]
&=
\frac{2}{3}P(A^c)
\end{align*}

Also:

\begin{align*}
P(A^c)
&=
\frac{0.6}{0.4}P(B^c) \\[2mm]
&=
\frac{3}{2}P(B^c)
\end{align*}

\textbf{4. Use the union formula to find $P(A)$}

The union formula is:

\begin{align*}
P(A^c\cup B^c)
&=
P(A^c)+P(B^c)-P(A^c\cap B^c)
\end{align*}

Use

\[
P(B^c)=\frac{2}{3}P(A^c)
\]

and

\[
P(A^c\cap B^c)=0.4P(A^c).
\]

Then:

\begin{align*}
0.8
&=
P(A^c)
+\frac{2}{3}P(A^c)
-0.4P(A^c) \\[2mm]
0.8
&=
P(A^c)
\left(
1+\frac{2}{3}-0.4
\right)
\end{align*}

Since $0.4=\frac{2}{5}$:

\begin{align*}
1+\frac{2}{3}-\frac{2}{5}
&=
\frac{15}{15}
+\frac{10}{15}
-\frac{6}{15} \\[2mm]
&=
\frac{19}{15}
\end{align*}

So:

\begin{align*}
0.8
&=
P(A^c)\frac{19}{15} \\[2mm]
P(A^c)
&=
\frac{0.8}{19/15} \\[2mm]
&=
\frac{12}{19}
\end{align*}

Therefore:

\begin{align*}
P(A)
&=
1-P(A^c) \\[2mm]
&=
1-\frac{12}{19} \\[2mm]
&=
\frac{7}{19} \\[2mm]
&\approx 0.37
\end{align*}

\textbf{5. Use the union formula to find $P(B)$}

Now use the other replacement:

\[
P(A^c)=\frac{3}{2}P(B^c)
\]

and

\[
P(A^c\cap B^c)=0.6P(B^c).
\]

Then:

\begin{align*}
0.8
&=
\frac{3}{2}P(B^c)
+P(B^c)
-0.6P(B^c) \\[2mm]
0.8
&=
P(B^c)
\left(
\frac{3}{2}+1-0.6
\right)
\end{align*}

The factor is:

\begin{align*}
\frac{3}{2}+1-0.6
&=
1.5+1-0.6 \\[2mm]
&=1.9
\end{align*}

So:

\begin{align*}
0.8
&=
1.9P(B^c) \\[2mm]
P(B^c)
&=
\frac{0.8}{1.9} \\[2mm]
&=
\frac{8}{19}
\end{align*}

Therefore:

\begin{align*}
P(B)
&=
1-P(B^c) \\[2mm]
&=
1-\frac{8}{19} \\[2mm]
&=
\frac{11}{19} \\[2mm]
&\approx 0.58
\end{align*}

\subsubsection*{Final answer}
\[
P(A)=\frac{7}{19}\approx 0.37.
\]

\[
P(B)=\frac{11}{19}\approx 0.58.
\]

\subsubsection*{What to remember}
Conditional probability:

\[
P(E\mid F)
=
\frac{P(E\cap F)}{P(F)}.
\]

Union formula:

\[
P(E\cup F)
=
P(E)+P(F)-P(E\cap F).
\]

Here the key trick is that

\[
A^c\cap B^c
=
B^c\cap A^c
\]

is the same overlap, so the two conditional formulas can be connected.

\section*{Uppgift 2}

\subsection*{Original question}
Betrakta funktionen

\[
f(x)=
\begin{cases}
ax^2,
& 0<x<3,\\
0,
& \text{annars.}
\end{cases}
\]

\begin{enumerate}[label=\alph*)]
\item
Bestäm $a$ så att $f$ bildar en täthetsfunktion för en s.v. $X$.
\hfill (1.5p)

\item
Bestäm fördelningsfunktionen $F_X$.
\hfill (1.5p)

\item
Beräkna $p(X>1)$.
\hfill (1p)

\item
Beräkna väntevärde och standardavvikelse för $X$.
\hfill (2p)
\end{enumerate}

\subsection*{Compiled notes from Yacoub + Obid}

\subsubsection*{What we are doing}
The notes say that $f$ is a density function, or PDF.
A PDF represents the density, and the total area under the graph
must be $1$.

First we choose $a$ so that the total area is $1$.
Then we integrate the PDF from the lower limit up to $x$ to get
the CDF.
After that we use the CDF to calculate $P(X>1)$.
Finally we calculate the expected value and standard deviation.

\subsubsection*{Given information}
The function is:

\[
f(x)=
\begin{cases}
ax^2,
& 0<x<3,\\
0,
& \text{otherwise.}
\end{cases}
\]

The interval where the density is nonzero is:

\[
0<x<3.
\]

\subsubsection*{Step-by-step solution}

\textbf{a) Find $a$}

For a PDF, the total area must be $1$:

\[
\int_0^3 ax^2\,dx=1.
\]

Integrate:

\begin{align*}
\int_0^3 ax^2\,dx
&=
\left[
a\frac{x^3}{3}
\right]_0^3 \\[2mm]
&=1
\end{align*}

Plug in the limits:

\begin{align*}
a\frac{3^3}{3}-0
&=1 \\[2mm]
a\frac{27}{3}
&=1 \\[2mm]
9a
&=1 \\[2mm]
a
&=\frac{1}{9}
\end{align*}

So the PDF is:

\[
f(x)=\frac{1}{9}x^2,
\qquad 0<x<3.
\]

\subsection*{Original question repeated for this part}
Betrakta funktionen

\[
f(x)=
\begin{cases}
ax^2,
& 0<x<3,\\
0,
& \text{annars.}
\end{cases}
\]

\begin{enumerate}[label=\alph*)]
\item
Bestäm $a$ så att $f$ bildar en täthetsfunktion för en s.v. $X$.
\hfill (1.5p)

\item
Bestäm fördelningsfunktionen $F_X$.
\hfill (1.5p)

\item
Beräkna $p(X>1)$.
\hfill (1p)

\item
Beräkna väntevärde och standardavvikelse för $X$.
\hfill (2p)
\end{enumerate}

\textbf{b) Find the CDF}

The notes say the CDF is the probability up to a point $x$.
The PDF represents the whole graph, while the CDF accumulates
the area up to $x$.

For $0<x<3$:

\begin{align*}
F_X(x)
&=
\int_0^x
\frac{1}{9}t^2\,dt \\[2mm]
&=
\frac{1}{9}
\left[
\frac{t^3}{3}
\right]_0^x \\[2mm]
&=
\frac{1}{9}
\left(
\frac{x^3}{3}-0
\right) \\[2mm]
&=
\frac{x^3}{27}
\end{align*}

So:

\[
F_X(x)=
\begin{cases}
0,
& x\leq 0,\\
\dfrac{x^3}{27},
& 0<x<3,\\
1,
& x\geq 3.
\end{cases}
\]

\subsection*{Original question repeated for this part}
Betrakta funktionen

\[
f(x)=
\begin{cases}
ax^2,
& 0<x<3,\\
0,
& \text{annars.}
\end{cases}
\]

\begin{enumerate}[label=\alph*)]
\item
Bestäm $a$ så att $f$ bildar en täthetsfunktion för en s.v. $X$.
\hfill (1.5p)

\item
Bestäm fördelningsfunktionen $F_X$.
\hfill (1.5p)

\item
Beräkna $p(X>1)$.
\hfill (1p)

\item
Beräkna väntevärde och standardavvikelse för $X$.
\hfill (2p)
\end{enumerate}

\textbf{c) Calculate $P(X>1)$}

The notes explain that $P(X>1)$ means the chance for the outcome
of the random variable to be bigger than $1$.

Use the CDF:

\begin{align*}
P(X>1)
&=
1-P(X\leq 1) \\[2mm]
&=
1-F_X(1)
\end{align*}

From part b:

\[
F_X(x)=\frac{x^3}{27}.
\]

Plug in $x=1$:

\begin{align*}
F_X(1)
&=
\frac{1^3}{27} \\[2mm]
&=
\frac{1}{27}
\end{align*}

Therefore:

\begin{align*}
P(X>1)
&=
1-\frac{1}{27} \\[2mm]
&=
\frac{26}{27} \\[2mm]
&\approx 0.963
\end{align*}

So the chance that the random variable gives an outcome bigger
than $1$ is about $96.3\,\%$.

\subsection*{Original question repeated for this part}
Betrakta funktionen

\[
f(x)=
\begin{cases}
ax^2,
& 0<x<3,\\
0,
& \text{annars.}
\end{cases}
\]

\begin{enumerate}[label=\alph*)]
\item
Bestäm $a$ så att $f$ bildar en täthetsfunktion för en s.v. $X$.
\hfill (1.5p)

\item
Bestäm fördelningsfunktionen $F_X$.
\hfill (1.5p)

\item
Beräkna $p(X>1)$.
\hfill (1p)

\item
Beräkna väntevärde och standardavvikelse för $X$.
\hfill (2p)
\end{enumerate}

\textbf{d) Expected value and standard deviation}

The notes describe expected value as the center, or mean:

\[
\mu=E[X].
\]

Standard deviation tells how much the values vary from the mean.

\textbf{Find $E[X]$}

The notes say to multiply $x$ with the PDF:

\begin{align*}
E[X]
&=
\int_0^3
x\left(\frac{1}{9}x^2\right)\,dx \\[2mm]
&=
\int_0^3
\frac{1}{9}x^3\,dx \\[2mm]
&=
\frac{1}{9}
\left[
\frac{x^4}{4}
\right]_0^3 \\[2mm]
&=
\frac{1}{9}
\left(
\frac{81}{4}-0
\right) \\[2mm]
&=
\frac{81}{36} \\[2mm]
&=
\frac{9}{4} \\[2mm]
&=2.25
\end{align*}

So:

\[
E[X]=2.25.
\]

\textbf{Find $E[X^2]$}

For standard deviation, first use:

\[
\operatorname{Var}(X)
=
E[X^2]-(E[X])^2.
\]

The notes say to multiply $x^2$ with the PDF:

\begin{align*}
E[X^2]
&=
\int_0^3
x^2\left(\frac{1}{9}x^2\right)\,dx \\[2mm]
&=
\int_0^3
\frac{1}{9}x^4\,dx \\[2mm]
&=
\frac{1}{9}
\left[
\frac{x^5}{5}
\right]_0^3 \\[2mm]
&=
\frac{1}{9}
\left(
\frac{243}{5}-0
\right) \\[2mm]
&=
5.4
\end{align*}

Now calculate the variance:

\begin{align*}
\operatorname{Var}(X)
&=
E[X^2]-(E[X])^2 \\[2mm]
&=
5.4-(2.25)^2 \\[2mm]
&=
5.4-5.0625 \\[2mm]
&=
0.3375
\end{align*}

Then:

\begin{align*}
\sigma
&=
\sqrt{\operatorname{Var}(X)} \\[2mm]
&=
\sqrt{0.3375} \\[2mm]
&\approx 0.581
\end{align*}

\subsubsection*{Final answer}
\begin{enumerate}[label=\alph*)]
\item
\[
a=\frac{1}{9}.
\]

\item
\[
F_X(x)=
\begin{cases}
0,
& x\leq 0,\\
\dfrac{x^3}{27},
& 0<x<3,\\
1,
& x\geq 3.
\end{cases}
\]

\item
\[
P(X>1)=\frac{26}{27}\approx 0.963.
\]

\item
\[
E[X]=2.25,
\qquad
\sigma\approx 0.581.
\]
\end{enumerate}

\subsubsection*{What to remember}
For a density function:

\[
\int f(x)\,dx=1.
\]

For a CDF:

\[
F_X(x)=P(X\leq x).
\]

For expected value:

\[
E[X]=\int x f(x)\,dx.
\]

For standard deviation:

\[
\sigma=
\sqrt{E[X^2]-(E[X])^2}.
\]

\section*{Uppgift 3}

\subsection*{Original question}
Låt $X_1,X_2\in \operatorname{P}(\mu)$, där $\mu=2$ och antag att
$X_1,X_2$ är oberoende s.v.

\begin{enumerate}[label=\alph*)]
\item
Beräkna $P(X_1=1)$ och $P(1\leq X_1<4)$.

\item
Beräkna $P(\max(X_1,X_2)\leq 2)$.

\item
Låt $Y=X_1-2X_2$. Beräkna $E(Y)$ och $V(Y)$.
\end{enumerate}

\subsection*{Compiled notes from Yacoub + Obid}

\subsubsection*{What we are doing}
The notes use the Poisson distribution.
They describe it as a tool to predict the odds of random
fluctuations.

If we know the mean, the Poisson formula gives the chance that
an exact event happens $k$ times in a specific window.

The notes also remind us that for a Poisson random variable:

\[
E[X]=\mu,
\qquad
\operatorname{Var}(X)=\mu.
\]

Here:

\[
\mu=2.
\]

\subsubsection*{Given information}
\[
X_1,X_2\in \operatorname{P}(2)
\]

\[
X_1 \text{ and } X_2
\text{ are independent.}
\]

The Poisson formula is:

\[
P(X=k)=
\frac{\mu^k e^{-\mu}}{k!}.
\]

With $\mu=2$:

\[
P(X=k)=
\frac{2^k e^{-2}}{k!}.
\]

\subsubsection*{Step-by-step solution}

\textbf{a) Calculate $P(X_1=1)$}

Use the formula with $k=1$:

\begin{align*}
P(X_1=1)
&=
\frac{2^1 e^{-2}}{1!} \\[2mm]
&=
2e^{-2} \\[2mm]
&\approx 0.2707
\end{align*}

\textbf{Calculate $P(1\leq X_1<4)$}

Since $X_1$ is discrete:

\[
1\leq X_1<4
\]

means:

\[
X_1=1,2,3.
\]

So:

\begin{align*}
P(1\leq X_1<4)
&=
P(X_1=1)
+P(X_1=2)
+P(X_1=3)
\end{align*}

The notes write:

\[
P(X_1=1)\approx 0.2707.
\]

For $X_1=2$:

\begin{align*}
P(X_1=2)
&=
\frac{2^2e^{-2}}{2!} \\[2mm]
&=
2e^{-2} \\[2mm]
&\approx 0.2707
\end{align*}

For $X_1=3$:

\begin{align*}
P(X_1=3)
&=
\frac{2^3e^{-2}}{3!} \\[2mm]
&=
\frac{8}{6}e^{-2} \\[2mm]
&\approx 0.1804
\end{align*}

Therefore:

\begin{align*}
P(1\leq X_1<4)
&\approx
0.2707+0.2707+0.1804 \\[2mm]
&=0.7218
\end{align*}

\subsection*{Original question repeated for this part}
Låt $X_1,X_2\in \operatorname{P}(\mu)$, där $\mu=2$ och antag att
$X_1,X_2$ är oberoende s.v.

\begin{enumerate}[label=\alph*)]
\item
Beräkna $P(X_1=1)$ och $P(1\leq X_1<4)$.

\item
Beräkna $P(\max(X_1,X_2)\leq 2)$.

\item
Låt $Y=X_1-2X_2$. Beräkna $E(Y)$ och $V(Y)$.
\end{enumerate}

\textbf{b) Calculate $P(\max(X_1,X_2)\leq 2)$}

The notes explain this as:
both $X_1$ and $X_2$ must be at most $2$.

So:

\[
\max(X_1,X_2)\leq 2
\]

means:

\[
X_1\leq 2
\quad \text{and}\quad
X_2\leq 2.
\]

Because $X_1$ and $X_2$ are independent:

\begin{align*}
P(\max(X_1,X_2)\leq 2)
&=
P(X_1\leq 2)\cdot P(X_2\leq 2)
\end{align*}

Find $P(X\leq 2)$:

\begin{align*}
P(X\leq 2)
&=
P(X=0)+P(X=1)+P(X=2)
\end{align*}

The notes calculate:

\begin{align*}
P(X=0)
&=
\frac{2^0e^{-2}}{0!}
=e^{-2} \\[2mm]
P(X=1)
&=
2e^{-2} \\[2mm]
P(X=2)
&=
2e^{-2}
\end{align*}

So:

\begin{align*}
P(X\leq 2)
&=
e^{-2}+2e^{-2}+2e^{-2} \\[2mm]
&=
5e^{-2}
\end{align*}

Thus:

\begin{align*}
P(\max(X_1,X_2)\leq 2)
&=
5e^{-2}\cdot 5e^{-2} \\[2mm]
&=
25e^{-4} \\[2mm]
&\approx 0.4578
\end{align*}

\subsection*{Original question repeated for this part}
Låt $X_1,X_2\in \operatorname{P}(\mu)$, där $\mu=2$ och antag att
$X_1,X_2$ är oberoende s.v.

\begin{enumerate}[label=\alph*)]
\item
Beräkna $P(X_1=1)$ och $P(1\leq X_1<4)$.

\item
Beräkna $P(\max(X_1,X_2)\leq 2)$.

\item
Låt $Y=X_1-2X_2$. Beräkna $E(Y)$ och $V(Y)$.
\end{enumerate}

\textbf{c) Calculate $E(Y)$ and $V(Y)$}

The variable is:

\[
Y=X_1-2X_2.
\]

The notes split the plus/minus parts in the variable.

For expectation:

\[
E(aX_1+bX_2)
=
aE[X_1]+bE[X_2].
\]

So:

\begin{align*}
E(Y)
&=
E(X_1-2X_2) \\[2mm]
&=
E(X_1)-2E(X_2) \\[2mm]
&=
1\cdot 2-2\cdot 2 \\[2mm]
&=
2-4 \\[2mm]
&=-2
\end{align*}

For variance, since $X_1$ and $X_2$ are independent:

\[
\operatorname{Var}(aX_1+bX_2)
=
a^2\operatorname{Var}(X_1)
+b^2\operatorname{Var}(X_2).
\]

Use $\operatorname{Var}(X_1)=2$ and
$\operatorname{Var}(X_2)=2$:

\begin{align*}
V(Y)
&=
\operatorname{Var}(X_1-2X_2) \\[2mm]
&=
1^2\operatorname{Var}(X_1)
+(-2)^2\operatorname{Var}(X_2) \\[2mm]
&=
1^2\cdot 2+4\cdot 2 \\[2mm]
&=
2+8 \\[2mm]
&=10
\end{align*}

\subsubsection*{Final answer}
\begin{enumerate}[label=\alph*)]
\item
\[
P(X_1=1)=2e^{-2}\approx 0.2707.
\]

\[
P(1\leq X_1<4)\approx 0.7218.
\]

\item
\[
P(\max(X_1,X_2)\leq 2)
=25e^{-4}\approx 0.4578.
\]

\item
\[
E(Y)=-2,
\qquad
V(Y)=10.
\]
\end{enumerate}

\subsubsection*{What to remember}
For Poisson:

\[
P(X=k)=
\frac{\mu^k e^{-\mu}}{k!}.
\]

For Poisson:

\[
E[X]=\mu,
\qquad
\operatorname{Var}(X)=\mu.
\]

For independent variables:

\[
P(X_1\leq 2,\ X_2\leq 2)
=
P(X_1\leq 2)P(X_2\leq 2).
\]

For variance of a linear combination:

\[
\operatorname{Var}(aX_1+bX_2)
=
a^2\operatorname{Var}(X_1)
+b^2\operatorname{Var}(X_2).
\]

\section*{Uppgift 4}

\subsection*{Original question}
Låt $x_1,x_2,\ldots,x_n$, där $0<x_i<1$ för alla $i$, vara ett
stickprov från en fördelning med täthetsfunktion

\[
f_X(x)=
\begin{cases}
\theta x^{\theta-1},
& 0<x<1,\\
0,
& \text{annars.}
\end{cases}
\]

Ange ML-skattningen av $\theta$.
\hfill (6p)

\subsection*{Compiled notes from Yacoub + Obid}

\subsubsection*{What we are doing}
The notes say we want to find the formula for $\theta$ that maximizes
the probability of getting exactly the data we observed.

This is maximum likelihood:

\begin{itemize}
\item build the likelihood function,
\item take the logarithm,
\item differentiate,
\item set the derivative equal to zero,
\item solve for $\theta$.
\end{itemize}

\subsubsection*{Given information}
The data are:

\[
x_1,x_2,\ldots,x_n
\]

and each value satisfies:

\[
0<x_i<1.
\]

The density is:

\[
f_X(x)=\theta x^{\theta-1},
\qquad 0<x<1.
\]

\subsubsection*{Step-by-step solution}

\textbf{Step 1: Likelihood function}

The notes use:

\[
L(\theta)
=
\prod_{i=1}^n f_X(x_i).
\]

Substitute the density:

\begin{align*}
L(\theta)
&=
\prod_{i=1}^n
\left(\theta x_i^{\theta-1}\right)
\end{align*}

Apply the product rule:

\begin{align*}
L(\theta)
&=
\left(\prod_{i=1}^n \theta\right)
\left(\prod_{i=1}^n x_i\right)^{\theta-1}
\end{align*}

Since there are $n$ copies of $\theta$:

\begin{align*}
L(\theta)
&=
\theta^n
\left(\prod_{i=1}^n x_i\right)^{\theta-1}
\end{align*}

\textbf{Step 2: Log-likelihood}

The notes take the logarithm:

\begin{align*}
\ln L(\theta)
&=
\ln\left[
\theta^n
\left(\prod_{i=1}^n x_i\right)^{\theta-1}
\right]
\end{align*}

Use:

\[
\ln(A\cdot B)=\ln(A)+\ln(B)
\]

and:

\[
\ln(A^b)=b\ln(A).
\]

Then:

\begin{align*}
\ln L(\theta)
&=
\ln(\theta^n)
+\ln\left[
\left(\prod_{i=1}^n x_i\right)^{\theta-1}
\right] \\[2mm]
&=
n\ln\theta
+(\theta-1)
\ln\left(\prod_{i=1}^n x_i\right)
\end{align*}

The notes also use:

\[
\ln\left(\prod x_i\right)
=
\sum \ln(x_i).
\]

So:

\begin{align*}
\ln L(\theta)
&=
n\ln\theta
+(\theta-1)
\sum_{i=1}^n \ln(x_i)
\end{align*}

\textbf{Step 3: Differentiate}

To find the maximum, take the derivative and set it equal to zero:

\begin{align*}
\frac{d}{d\theta}\ln L(\theta)
&=
\frac{d}{d\theta}
\left[
n\ln\theta
+(\theta-1)
\sum_{i=1}^n \ln(x_i)
\right]
\end{align*}

Differentiate:

\begin{align*}
\frac{d}{d\theta}\ln L(\theta)
&=
\frac{n}{\theta}
+1\cdot
\sum_{i=1}^n \ln(x_i)
\end{align*}

\textbf{Step 4: Set equal to zero and solve}

\begin{align*}
\frac{n}{\theta}
+\sum_{i=1}^n \ln(x_i)
&=0
\end{align*}

Move the sum to the other side:

\begin{align*}
\frac{n}{\theta}
&=
-\sum_{i=1}^n \ln(x_i)
\end{align*}

Solve for $\theta$:

\begin{align*}
\theta
&=
\frac{n}
{-\sum_{i=1}^n \ln(x_i)}
\end{align*}

\subsubsection*{Final answer}
\[
\hat{\theta}_{\mathrm{ML}}
=
\frac{n}
{-\sum_{i=1}^n \ln(x_i)}.
\]

\subsubsection*{What to remember}
For ML estimation:

\[
L(\theta)=
\prod_{i=1}^n f_X(x_i).
\]

Taking logs turns products into sums:

\[
\ln\left(\prod x_i\right)
=
\sum \ln(x_i).
\]

To maximize, differentiate the log-likelihood and set it equal to zero.

\section*{Uppgift 5}

\subsection*{Original question}
Följande 8 observationer kan antas vara observationer från
$N(\mu,\sigma)$:

\[
13.4,\;14.1,\;15.1,\;14.3,\;14.9,\;13.9,\;14.4,\;14.0.
\]

Beräkning ger att

\[
\sum_{i=1}^{8} x_i=114.1,
\qquad
\sum_{i=1}^{8} x_i^2=1629.45.
\]

\begin{enumerate}[label=\alph*)]
\item
Ange en lämplig punktskattning av $\sigma^2$ och $\sigma$.
\hfill (2p)

\item
Ange ett $90\,\%$ konfidensintervall för $\sigma$.
\hfill (4p)
\end{enumerate}

\subsection*{Compiled notes from Yacoub + Obid}

\subsubsection*{What we are doing}
The notes use descriptive statistics for a normal sample.
First we estimate the variance and standard deviation from the sample.

Then we build a $90\,\%$ confidence interval for $\sigma$.
The notes first find the interval for $\sigma^2$ using the
chi-square distribution, and then take square roots to get the
interval for $\sigma$.

\subsubsection*{Given information}
From the question:

\[
n=8
\]

\[
\sum x_i=114.1
\]

\[
\sum x_i^2=1629.45
\]

For part b:

\[
90\,\%\text{ confidence}
\]

so

\[
\alpha=0.10.
\]

\subsubsection*{Step-by-step solution}

\textbf{a) Point estimate of $\sigma^2$ and $\sigma$}

The notes use the sample variance:

\[
s^2
=
\frac{1}{n-1}
\left(
\sum x_i^2-n\bar{x}^{\,2}
\right).
\]

\textbf{Step 1: Find the sample mean}

\begin{align*}
\bar{x}
&=
\frac{\sum x_i}{n} \\[2mm]
&=
\frac{114.1}{8} \\[2mm]
&=
14.2625 \\[2mm]
&\approx 14.26
\end{align*}

\textbf{Step 2: Calculate the numerator}

\begin{align*}
\sum x_i^2-n\bar{x}^{\,2}
&=
1629.45
-8(14.2625)^2 \\[2mm]
&=
2.09875
\end{align*}

\textbf{Step 3: Divide by $n-1$}

\begin{align*}
s^2
&=
\frac{2.09875}{8-1} \\[2mm]
&=
\frac{2.09875}{7} \\[2mm]
&\approx 0.2998
\end{align*}

So the point estimate of $\sigma^2$ is:

\[
s^2\approx 0.2998.
\]

For $\sigma$, take the square root:

\begin{align*}
s
&=
\sqrt{0.2998} \\[2mm]
&\approx 0.5476
\end{align*}

So the point estimate of $\sigma$ is:

\[
s\approx 0.5476.
\]

\subsection*{Original question repeated for this part}
Följande 8 observationer kan antas vara observationer från
$N(\mu,\sigma)$:

\[
13.4,\;14.1,\;15.1,\;14.3,\;14.9,\;13.9,\;14.4,\;14.0.
\]

Beräkning ger att

\[
\sum_{i=1}^{8} x_i=114.1,
\qquad
\sum_{i=1}^{8} x_i^2=1629.45.
\]

\begin{enumerate}[label=\alph*)]
\item
Ange en lämplig punktskattning av $\sigma^2$ och $\sigma$.
\hfill (2p)

\item
Ange ett $90\,\%$ konfidensintervall för $\sigma$.
\hfill (4p)
\end{enumerate}

\textbf{b) 90\% confidence interval for $\sigma$}

The notes use the chi-square confidence interval for variance:

\[
\left[
\frac{(n-1)s^2}{\chi^2_{\text{upper}}},
\frac{(n-1)s^2}{\chi^2_{\text{lower}}}
\right].
\]

Here:

\[
df=n-1=8-1=7.
\]

For a $90\,\%$ confidence interval:

\[
\alpha=0.10.
\]

Split the error into two tails:

\[
0.10/2=0.05.
\]

From the chi-square table with $df=7$, the notes use:

\[
\chi^2_{0.95}=14.067
\]

\[
\chi^2_{0.05}=2.167
\]

Also, from part a:

\[
(n-1)s^2=2.09875.
\]

\textbf{Lower limit for $\sigma^2$}

\begin{align*}
\frac{(n-1)s^2}{\chi^2_{0.95}}
&=
\frac{2.09875}{14.067} \\[2mm]
&\approx 0.1492
\end{align*}

\textbf{Upper limit for $\sigma^2$}

\begin{align*}
\frac{(n-1)s^2}{\chi^2_{0.05}}
&=
\frac{2.09875}{2.167} \\[2mm]
&\approx 0.9685
\end{align*}

So the confidence interval for the variance is:

\[
\sigma^2\in [0.1492,\;0.9685].
\]

To get the interval for $\sigma$, take square roots:

\begin{align*}
\sigma
&\in
\left[
\sqrt{0.1492},
\sqrt{0.9685}
\right] \\[2mm]
&\approx
[0.386,\;0.984].
\end{align*}

\subsubsection*{Final answer}
\begin{enumerate}[label=\alph*)]
\item
\[
\widehat{\sigma^2}=s^2\approx 0.2998.
\]

\[
\hat{\sigma}=s\approx 0.5476.
\]

\item
\[
90\,\%\text{ confidence interval for }\sigma:
\]

\[
\sigma\in[0.386,\;0.984].
\]
\end{enumerate}

\subsubsection*{What to remember}
Sample variance:

\[
s^2
=
\frac{1}{n-1}
\left(
\sum x_i^2-n\bar{x}^{\,2}
\right).
\]

For a normal sample, a confidence interval for $\sigma^2$ uses
the chi-square distribution:

\[
\left[
\frac{(n-1)s^2}{\chi^2_{1-\alpha/2}},
\frac{(n-1)s^2}{\chi^2_{\alpha/2}}
\right].
\]

For $\sigma$, take the square root of both limits.

\section*{Uppgift 6}

\subsection*{Original question}
Vikten av en flygpassagerare (inklusive handbagage) antas vara
normalfördelad $N(\mu,\sigma)$. Under många år antogs $\mu=80$ kg,
men man vill nu testa om medelvikten har ökat.
Vi kan anta att $\sigma=10$.

\begin{enumerate}[label=\alph*)]
\item
Undersök om $\mu>80$ genom att ställa upp ett lämpligt hypotestest,
med signifikansnivån $5\,\%$ (nollhypotes och mothypotes ska vara
tydligt formulerade).
\hfill (2p)

\item
Utför ditt hypotestest i a) då $n=100$ och $\bar{x}=83$.
\hfill (2p)

\item
Antag att man tror att medelvikten ökat med 2 kg
(alltså $\mu=82$). Hur stort måste $n$ väljas för att kunna
upptäcka denna ökning med $80\,\%$ sannolikhet?
Alltså, hur stort måste $n$ väljas så att
$P(\text{Förkasta }H_0)>0.8$, under antagandet att $\mu=82$?
(Du kan här använda att $\Phi(0.84)\approx 0.8$.)
\hfill (2p)
\end{enumerate}

\subsection*{Compiled notes from Yacoub + Obid}

\subsubsection*{What we are doing}
The notes set up a hypothesis test for the mean weight.
The old mean is $80$ kg, and we test if the mean has increased.

Since $\sigma=10$ is known, the notes use a normal $z$ test.

Part c is about statistical power:
we choose $n$ large enough so that if the true mean is $82$,
we reject $H_0$ with probability at least $80\,\%$.

\subsubsection*{Given information}
\[
X\sim N(\mu,\sigma)
\]

\[
\sigma=10
\]

Old assumed mean:

\[
\mu=80
\]

Significance level:

\[
\alpha=5\,\%=0.05
\]

\subsubsection*{Step-by-step solution}

\textbf{a) Set up the hypothesis test}

The notes write:

\[
H_0:\mu=80
\]

\[
H_1:\mu>80
\]

This is a one-sided test, because the question asks if the mean
weight has increased.

The significance level is:

\[
\alpha=0.05.
\]

The notes explain this as:
$H_0$ is rejected only when there is less than a $5\,\%$ chance that
the result happened by luck.

\subsection*{Original question repeated for this part}
Vikten av en flygpassagerare (inklusive handbagage) antas vara
normalfördelad $N(\mu,\sigma)$. Under många år antogs $\mu=80$ kg,
men man vill nu testa om medelvikten har ökat.
Vi kan anta att $\sigma=10$.

\begin{enumerate}[label=\alph*)]
\item
Undersök om $\mu>80$ genom att ställa upp ett lämpligt hypotestest,
med signifikansnivån $5\,\%$ (nollhypotes och mothypotes ska vara
tydligt formulerade).
\hfill (2p)

\item
Utför ditt hypotestest i a) då $n=100$ och $\bar{x}=83$.
\hfill (2p)

\item
Antag att man tror att medelvikten ökat med 2 kg
(alltså $\mu=82$). Hur stort måste $n$ väljas för att kunna
upptäcka denna ökning med $80\,\%$ sannolikhet?
Alltså, hur stort måste $n$ väljas så att
$P(\text{Förkasta }H_0)>0.8$, under antagandet att $\mu=82$?
(Du kan här använda att $\Phi(0.84)\approx 0.8$.)
\hfill (2p)
\end{enumerate}

\textbf{b) Perform the test for $n=100$ and $\bar{x}=83$}

Given:

\[
n=100,
\qquad
\bar{x}=83.
\]

The standard error is:

\begin{align*}
SE
&=
\frac{\sigma}{\sqrt{n}} \\[2mm]
&=
\frac{10}{\sqrt{100}} \\[2mm]
&=
\frac{10}{10} \\[2mm]
&=1
\end{align*}

The $z$-score is:

\begin{align*}
z
&=
\frac{\bar{x}-\mu}{SE} \\[2mm]
&=
\frac{83-80}{1} \\[2mm]
&=3
\end{align*}

For a one-sided $5\,\%$ test, the notes use the normal table:

\[
z_{0.95}\approx 1.645.
\]

Compare:

\[
3>1.645.
\]

So we reject $H_0$.

There is evidence that the average weight has increased.

\subsection*{Original question repeated for this part}
Vikten av en flygpassagerare (inklusive handbagage) antas vara
normalfördelad $N(\mu,\sigma)$. Under många år antogs $\mu=80$ kg,
men man vill nu testa om medelvikten har ökat.
Vi kan anta att $\sigma=10$.

\begin{enumerate}[label=\alph*)]
\item
Undersök om $\mu>80$ genom att ställa upp ett lämpligt hypotestest,
med signifikansnivån $5\,\%$ (nollhypotes och mothypotes ska vara
tydligt formulerade).
\hfill (2p)

\item
Utför ditt hypotestest i a) då $n=100$ och $\bar{x}=83$.
\hfill (2p)

\item
Antag att man tror att medelvikten ökat med 2 kg
(alltså $\mu=82$). Hur stort måste $n$ väljas för att kunna
upptäcka denna ökning med $80\,\%$ sannolikhet?
Alltså, hur stort måste $n$ väljas så att
$P(\text{Förkasta }H_0)>0.8$, under antagandet att $\mu=82$?
(Du kan här använda att $\Phi(0.84)\approx 0.8$.)
\hfill (2p)
\end{enumerate}

\textbf{c) Required sample size}

The notes say this is about statistical power.
We want $n$ large enough so that, if the true mean is $82$,
we catch it $80\,\%$ of the time.

Given:

\[
\mu=82
\]

\[
P(\text{reject }H_0)>0.8
\]

and

\[
\Phi(0.84)\approx 0.8.
\]

The rejection rule from the one-sided test is:

\[
\bar{x}>
80+1.645\cdot \frac{10}{\sqrt{n}}.
\]

The notes write the power requirement as:

\[
P\left(
\bar{x}>
80+1.645\cdot \frac{10}{\sqrt{n}}
\right)
=0.8,
\qquad \text{given }\mu=82.
\]

Standardize using $\mu=82$:

\[
P\left(
Z>
\frac{
\left(80+1.645\cdot \frac{10}{\sqrt{n}}\right)-82
}
{10/\sqrt{n}}
\right)
=0.8.
\]

This becomes:

\[
P\left(
Z>
\frac{
-2+1.645\cdot \frac{10}{\sqrt{n}}
}
{10/\sqrt{n}}
\right)
=0.8.
\]

\textbf{English clarification:}
The notes stop here, but the next algebra step is directly implied.
Since $P(Z>-0.84)\approx 0.8$, set the inside value equal to $-0.84$.

\begin{align*}
\frac{
-2+1.645\cdot \frac{10}{\sqrt{n}}
}
{10/\sqrt{n}}
&=-0.84
\end{align*}

Split the fraction:

\begin{align*}
\frac{-2}{10/\sqrt{n}}
+1.645
&=-0.84 \\[2mm]
-\frac{2\sqrt{n}}{10}
+1.645
&=-0.84
\end{align*}

Move terms:

\begin{align*}
-\frac{2\sqrt{n}}{10}
&=
-0.84-1.645 \\[2mm]
-\frac{2\sqrt{n}}{10}
&=
-2.485
\end{align*}

So:

\begin{align*}
\frac{2\sqrt{n}}{10}
&=2.485 \\[2mm]
\sqrt{n}
&=
\frac{2.485\cdot 10}{2} \\[2mm]
&=12.425
\end{align*}

Square both sides:

\begin{align*}
n
&=
(12.425)^2 \\[2mm]
&\approx 154.38
\end{align*}

Since $n$ must be a whole number:

\[
n=155.
\]

\subsubsection*{Final answer}
\begin{enumerate}[label=\alph*)]
\item
\[
H_0:\mu=80,
\qquad
H_1:\mu>80.
\]

\item
\[
z=3.
\]
Since $3>1.645$, reject $H_0$.
There is evidence that the mean weight has increased.

\item
\[
n\geq 155.
\]
So choose
\[
n=155.
\]
\end{enumerate}

\subsubsection*{What to remember}
For a one-sided test with known $\sigma$:

\[
z=
\frac{\bar{x}-\mu_0}{\sigma/\sqrt{n}}.
\]

At significance level $5\,\%$ for an upper-tail test:

\[
z_{0.95}\approx 1.645.
\]

For power calculations, write the rejection rule in terms of
$\bar{x}$, then standardize using the alternative mean.

\end{document}
